\documentclass[a4paper,10pt]{article}
\usepackage{sk-techreport}

% Diagrams
% Shared TikZ styles for sk-techreport diagrams.
% Vendored into this report directory for portability.
%
% Included via % Shared TikZ styles for sk-techreport diagrams.
% Vendored into this report directory for portability.
%
% Included via % Shared TikZ styles for sk-techreport diagrams.
% Vendored into this report directory for portability.
%
% Included via \input{sk-diagrams.tex}. Do not use \usepackage here.

\tikzset{
  skd/.style={font=\sffamily},
  skdBox/.style={skd, draw=black!65, rounded corners=2pt, thick, inner sep=6pt, align=center},
  skdData/.style={skdBox, fill=layerE!35},
  skdCompute/.style={skdBox, fill=layerB!35},
  skdControl/.style={skdBox, fill=layerC!30},
  skdStore/.style={skdBox, fill=layerD!25},
  skdNote/.style={skd, draw=black!45, fill=black!3, rounded corners=2pt, inner sep=5pt, align=left},
  skdArrow/.style={-{Latex}, thick, draw=black!75},
  skdDataArrow/.style={skdArrow, draw=black!70},
  skdCtrlArrow/.style={skdArrow, dashed, draw=black!65},
  skdBackpressure/.style={skdArrow, draw=red!65!black},
  skdPhase/.style={skd, draw=black!35, fill=black!2, rounded corners=2pt, inner sep=3pt},
}
. Do not use \usepackage here.

\tikzset{
  skd/.style={font=\sffamily},
  skdBox/.style={skd, draw=black!65, rounded corners=2pt, thick, inner sep=6pt, align=center},
  skdData/.style={skdBox, fill=layerE!35},
  skdCompute/.style={skdBox, fill=layerB!35},
  skdControl/.style={skdBox, fill=layerC!30},
  skdStore/.style={skdBox, fill=layerD!25},
  skdNote/.style={skd, draw=black!45, fill=black!3, rounded corners=2pt, inner sep=5pt, align=left},
  skdArrow/.style={-{Latex}, thick, draw=black!75},
  skdDataArrow/.style={skdArrow, draw=black!70},
  skdCtrlArrow/.style={skdArrow, dashed, draw=black!65},
  skdBackpressure/.style={skdArrow, draw=red!65!black},
  skdPhase/.style={skd, draw=black!35, fill=black!2, rounded corners=2pt, inner sep=3pt},
}
. Do not use \usepackage here.

\tikzset{
  skd/.style={font=\sffamily},
  skdBox/.style={skd, draw=black!65, rounded corners=2pt, thick, inner sep=6pt, align=center},
  skdData/.style={skdBox, fill=layerE!35},
  skdCompute/.style={skdBox, fill=layerB!35},
  skdControl/.style={skdBox, fill=layerC!30},
  skdStore/.style={skdBox, fill=layerD!25},
  skdNote/.style={skd, draw=black!45, fill=black!3, rounded corners=2pt, inner sep=5pt, align=left},
  skdArrow/.style={-{Latex}, thick, draw=black!75},
  skdDataArrow/.style={skdArrow, draw=black!70},
  skdCtrlArrow/.style={skdArrow, dashed, draw=black!65},
  skdBackpressure/.style={skdArrow, draw=red!65!black},
  skdPhase/.style={skd, draw=black!35, fill=black!2, rounded corners=2pt, inner sep=3pt},
}


\addbibresource{cliscrape-techreport.bib}

\begin{document}

\maketechreporttitle{
  title       = {cliscrape: A Rust CLI Parsing Engine with Interactive TUI Debugger},
  subtitle    = {TextFSM-Compatible Parsing, Modern Templates, and Network-Compiler Foundations},
  reportid    = {CS-TR-2026-01},
  classification = {OPEN SOURCE},
  version     = {0.1.2 report revision},
  date        = {April 2026},
  author      = {Simon Knight},
  abstract    = {Network automation pipelines rely heavily on parsing
    unstructured CLI output from network devices.  This report presents
    cliscrape, a Rust implementation of a TextFSM-compatible state machine
    parser. Historical benchmarks have reported multi-million-line-per-second
    throughput, but release performance claims should be regenerated from the
    current benchmark suite before publication. The current implementation
    supports legacy TextFSM templates, a modern YAML/TOML template format,
    embedded/XDG template resolution, structured logging, policy constraints,
    semantic diffing, grammar induction, best-effort synthetic generation, and
    an interactive TUI debugger. Full NTC-Templates coverage, universal ledger
    schemas, strict oracle semantics, release automation, and streaming parsing
    remain tracked follow-up work.},
  coverimage  = {
    \begin{tikzpicture}[scale=1.0, every node/.style={transform shape}]
      \node[skNode, minimum width=3cm, fill=layerA] (input) at (0,0) {\textbf{CLI Output}};
      \node[skNode, minimum width=3cm, fill=layerB] (parser) at (0,-2) {\textbf{FSM Engine}};
      \node[skNode, minimum width=2cm, fill=layerC, font=\tiny\sffamily] (textfsm) at (-2,-4) {TextFSM};
      \node[skNode, minimum width=2cm, fill=layerC, font=\tiny\sffamily] (yaml) at (2,-4) {YAML/TOML};
      \node[skNode, minimum width=3cm, fill=layerD] (output) at (0,-6) {\textbf{Structured Data}};
      \node[skNode, minimum width=2.5cm, fill=layerE] (tui) at (4,-2) {\textbf{TUI Debugger}};

      \draw[skdArrow, <->] (input) -- (parser);
      \draw[arrow] (parser) -- (textfsm);
      \draw[arrow] (parser) -- (yaml);
      \draw[arrow] (textfsm) -- (output);
      \draw[arrow] (yaml) -- (output);
      \draw[arrow] (parser) -- (tui);
    \end{tikzpicture}
  },
  revisions = {
    0.1.0 & March 2026 & Initial architecture report \\
    0.1.1 & March 2026 & Added ecosystem context, case studies, production hardening, and performance data \\
    0.1.2 & April 2026 & Synced claims with current code, tests, and roadmap gaps \\
  },
}

%% Acronym definitions
\begin{acronym}[TextFSM]
  \acro{FSM}{Finite State Machine}
  \acro{CLI}{Command-Line Interface}
  \acro{TUI}{Terminal User Interface}
  \acro{NTC}{Network to Code}
  \acro{JSON}{JavaScript Object Notation}
  \acro{YAML}{YAML Ain't Markup Language}
  \acro{TOML}{Tom's Obvious, Minimal Language}
\end{acronym}

\clearpage

%% =====================================================================
\section{Introduction and Ecosystem Context}
\label{sec:intro}
%% =====================================================================

The networking industry relies on CLI-based management for device
configuration and monitoring.  Parsing this unstructured text output
is often the bottleneck in automation pipelines.
TextFSM~\cite{textfsm2024} provides a template-based approach using
state machines, and the NTC-Templates~\cite{ntc_templates2024} library
offers hundreds of community-maintained templates. \texttt{cliscrape}
implements a compatible TextFSM loading path and a modern template format;
full coverage of the external NTC template corpus is still a compatibility
target rather than the current embedded library state.

However, the Python TextFSM implementation has three core limitations:
performance limitations due to interpreted execution (often struggling
with large \texttt{show tech-support} outputs), readability issues caused
by dense inline regex syntax, and a lack of observability when templates
fail to match expected output.

\subsection{Ecosystem Role}

\texttt{cliscrape} operates within the Operations Path of the broader
network automation ecosystem. It sits directly behind the connectivity layer,
transforming raw device output into structured operational state.

\begin{figure}[H]
\centering
\adjustbox{max width=\linewidth}{%
\begin{tikzpicture}[skd, node distance=12mm and 15mm]
  \node[skdData, minimum width=3cm] (device) {\textbf{Live Device}\\\footnotesize network hardware};
  \node[skdCompute, right=of device, minimum width=3.5cm] (interact) {\textbf{deviceinteraction}\\\footnotesize connect + orchestrate};
  \node[skdCompute, right=of interact, minimum width=3.5cm] (scrape) {\textbf{cliscrape}\\\footnotesize parse + extract};
  \node[skdStore, right=of scrape, minimum width=3.5cm] (analyze) {\textbf{autonetkit / netassure}\\\footnotesize topology + verify};

  \draw[skdDataArrow] (device) -- node[above, font=\scriptsize] {CLI/SSH} (interact);
  \draw[skdDataArrow] (interact) -- node[above, font=\scriptsize] {raw text} (scrape);
  \draw[skdDataArrow] (scrape) -- node[above, font=\scriptsize] {structured \ac{JSON}} (analyze);
\end{tikzpicture}%
}
\caption{The Operations Path pipeline. \texttt{deviceinteraction} handles connectivity, while \texttt{cliscrape} handles pure parsing, feeding operational state to analysis tools.}
\label{fig:cs-ecosystem}
\end{figure}

%% =====================================================================
\section{Operational Model and Case Studies}
\label{sec:operations}
%% =====================================================================

The foundational goal of \texttt{cliscrape} is taking semi-structured CLI text and enforcing a typed schema contract. Consider a standard \texttt{show version} command from a Cisco IOS device:

\begin{lstlisting}[language=bash, title={Raw CLI Output (Cisco IOS)}]
Cisco IOS Software, C2960 Software (C2960-LANBASEK9-M), Version 15.0(2)SE11, RELEASE SOFTWARE (fc3)
ROM: Bootstrap program is C2960 boot loader
SW-CORE-01 uptime is 5 weeks, 2 days, 3 hours, 45 minutes
System image file is "flash:/c2960-lanbasek9-mz.150-2.SE11.bin"
\end{lstlisting}

Instead of brittle custom scripts or string-splitting, the \ac{FSM} engine applies an embedded YAML or TextFSM template. The variables defined in the template dictate the keys of the resulting \ac{JSON} object.

\begin{lstlisting}[language=json, title={Parsed Structured Output (excerpt)}]
[
  {
    "hostname": "SW-CORE-01",
    "uptime": "5 weeks, 2 days, 3 hours, 45 minutes",
    "version": "15.0(2)SE11,"
  }
]
\end{lstlisting}

This operational model allows downstream tools to perform type-aware assertions against properties like \texttt{uptime} and \texttt{version}, decoupled from many vendor-specific formatting quirks of the original CLI. The current validation snapshots have also exposed template-specific extraction bugs, such as Cisco IOS serial-number handling, which should be fixed as template correctness work rather than hidden in report examples.

%% =====================================================================
\section{FSM Engine Architecture}
\label{sec:engine}
%% =====================================================================

\begin{figure}[H]
\centering
\adjustbox{max width=\linewidth}{%
\begin{tikzpicture}[skd, node distance=10mm and 12mm]
  \node[skdData, minimum width=3.6cm] (script) {\textbf{Template / Spec}\\\footnotesize states + rules};
  \node[skdCompute, right=of script, minimum width=3.6cm] (fsm) {\textbf{FSM Runtime}\\\footnotesize current state\\match+action loop};
  \node[skdCompute, right=of fsm, minimum width=3.6cm] (input) {\textbf{CLI Input}\\\footnotesize line stream};

  \node[skdCompute, below=of fsm, minimum width=5.4cm] (parse) {\textbf{Rule Match}\\\footnotesize regex + captures};
  \node[skdStore, right=of parse, minimum width=5.4cm] (out) {\textbf{Records}\\\footnotesize results + fields};

  \draw[skdDataArrow] (script) -- (fsm);
  \draw[skdDataArrow] (input.south) |- (fsm.east);
  \draw[skdDataArrow] (fsm) -- (parse);
  \draw[skdDataArrow] (parse) -- (out);
  \draw[skdCtrlArrow] (parse.north) -- node[right, font=\scriptsize] {Next/Continue/Record} (fsm.south);

  \node[skdNote, below=10mm of out, text width=11.2cm] {\textbf{Control vs data.} The same regex match both extracts data (captures) and drives control (state transitions). This is what makes TextFSM robust to prompts, pagination, and error banners.};
\end{tikzpicture}%
}
\caption{FSM engine model: rule matching is both evidence extraction and control-flow.}
\label{fig:cs-fsm-arch}
\end{figure}

The core engine implements the exact TextFSM state machine semantics using Rust.

\begin{figure}[H]
\centering
\adjustbox{max width=\linewidth}{%
\begin{tikzpicture}[skd, node distance=9mm and 12mm]
  \node[skdPhase, minimum width=11.2cm] (hdr) {\textbf{TextFSM semantics as a state machine over (line index, record state)}};

  \node[skdStore, below left=of hdr, minimum width=3.6cm] (state) {\textbf{State}\\\footnotesize current\_state\\record vars};
  \node[skdData, right=of state, minimum width=3.6cm] (line) {\textbf{Input cursor}\\\footnotesize line $i$};
  \node[skdCompute, right=of line, minimum width=3.6cm] (rules) {\textbf{Rule scan}\\\footnotesize ordered match};

  \node[skdControl, below=of line, minimum width=11.2cm] (acts) {\textbf{Actions}\\\footnotesize Next: $i\!\leftarrow\!i+1$\quad Continue: try next rule\quad Record: emit + clear\quad Clear: reset vars};

  \draw[skdDataArrow] (state) -- (rules);
  \draw[skdDataArrow] (line) -- (rules);
  \draw[skdCtrlArrow] (rules) -- (acts);
  \draw[skdCtrlArrow] (acts.west) |- node[left, font=\scriptsize] {updates} (state.south);
  \draw[skdCtrlArrow] (acts.east) |- node[right, font=\scriptsize] {advances} (line.south);
\end{tikzpicture}%
}
\caption{TextFSM parsing as an interpreter over a line cursor and an accumulating record.}
\label{fig:cs-semantics}
\end{figure}

For each input line, the engine evaluates rules in the current state
sequentially. On match, captured groups are assigned to named
variables, and the defined action (Next, Continue, Record, Clear) is executed.

\subsection{Stateful Modifiers}

Simple regular expressions are insufficient for network CLI outputs that exhibit hierarchical or stateful formatting (e.g., a routing table where a single routing protocol applies to multiple subsequent subnet lines). To support this, the engine implements TextFSM's stateful variable modifiers:
\begin{itemize}
    \item \textbf{\texttt{Filldown}:} Retains a captured variable's value across multiple record emissions until explicitly cleared or overwritten.
    \item \textbf{\texttt{List}:} Rather than overwriting a variable on subsequent matches, accumulates the matched values into a native \ac{JSON} array.
\end{itemize}
These modifiers allow the parser to build deep, relational data structures from flat text without requiring complex post-processing.

\insight{The Rust \texttt{regex} crate's DFA-based engine provides
  predictable $O(n)$ matching per line, avoiding the catastrophic
  backtracking possible with Python's \texttt{re} module on
  pathological patterns.}

%% =====================================================================
\section{Template Formats}
\label{sec:templates}
%% =====================================================================

\begin{figure}[H]
\centering
\adjustbox{max width=\linewidth}{%
\begin{tikzpicture}[skd, node distance=9mm and 12mm]
  \node[skdData, minimum width=4.0cm] (tmpl) {\textbf{Template}\\\footnotesize YAML/TextFSM\\named captures};
  \node[skdCompute, right=of tmpl, minimum width=4.0cm] (match) {\textbf{Matcher}\\\footnotesize apply to output\\extract groups};
  \node[skdStore, right=of match, minimum width=3.2cm] (fields) {\textbf{Fields}\\\footnotesize typed keys};
  \node[skdStore, below=of match, minimum width=11.2cm] (schema) {\textbf{Schema Contract}\\\footnotesize stable field names enable diffs, regression tests, and downstream automation};

  \draw[skdDataArrow] (tmpl) -- (match);
  \draw[skdDataArrow] (match) -- (fields);
  \draw[skdCtrlArrow] (fields.south) |- (schema.east);
\end{tikzpicture}%
}
\caption{Templates define the parsing contract: semistructured CLI output becomes stable, typed records.}
\label{fig:cs-templates}
\end{figure}

Supporting both TextFSM and YAML formats allows teams to migrate incrementally. Existing templates work immediately via a compatibility translation layer. New templates can use the more readable YAML format. Both compile to the same internal \ac{FSM} representation (IR).

Beyond mere readability, the shift to YAML/TOML introduces explicit schema definition capability. Because values can be natively typed (e.g., explicitly defining an interface MTU as an integer rather than an arbitrary regex string), \texttt{cliscrape} can support downstream tooling such as JSON Schema generation, typed client bindings, and documentation validation. Those assets are roadmap items, not all current shipped behavior.

\subsection{Native Type Coercion}

Legacy Python TextFSM implementations fundamentally treat all captured data as strings. If a device outputs an MTU of \texttt{1500}, the resulting JSON output will be \texttt{"mtu": "1500"}. This pushes the burden of data normalisation onto downstream analysis tools. 

Because \texttt{cliscrape}'s modern templates allow explicit type definitions, the engine performs native type coercion during the match phase. An extracted string representing a number is safely converted and emitted as a native \ac{JSON} Integer. This seemingly small feature represents a major architectural win for the Operations Path: when data reaches \texttt{netassure}, it is instantly queryable without secondary casting operations.

\begin{figure}[H]
\centering
\adjustbox{max width=\linewidth}{%
\begin{tikzpicture}[skd, node distance=9mm and 12mm]
  \node[skdPhase, minimum width=11.2cm] (hdr) {\textbf{Template compilation: two syntaxes, one executable IR}};

  \node[skdData, below left=of hdr, minimum width=3.6cm] (tfsm) {\textbf{TextFSM DSL}\\\footnotesize legacy templates};
  \node[skdData, right=of tfsm, minimum width=3.6cm] (yaml) {\textbf{YAML/TOML}\\\footnotesize readable templates};
  \node[skdCompute, right=of yaml, minimum width=3.6cm] (ir) {\textbf{IR}\\\footnotesize states + rules\\named captures};
  \node[skdCompute, below=of ir, minimum width=11.2cm] (run) {\textbf{Runtime}\\\footnotesize deterministic match loop; stable schema contract};

  \draw[skdDataArrow] (tfsm) -- (ir);
  \draw[skdDataArrow] (yaml) -- (ir);
  \draw[skdDataArrow] (ir) -- (run);
\end{tikzpicture}%
}
\caption{Compilation boundary: both template syntaxes compile into one IR consumed by the same runtime.}
\label{fig:cs-compile}
\end{figure}


%% =====================================================================
\section{Production Hardening}
\label{sec:hardening}
%% =====================================================================

Moving from a parsing prototype to a robust operational tool required several critical system enhancements to handle edge cases, simplify deployments, and ensure continuous availability.

\subsection{Native Binary Deployment}
In enterprise networking environments, deploying tools to jump-hosts or restricted bastion servers is often hindered by Python runtime dependencies, virtual environments, and missing \texttt{pip} packages. Because \texttt{cliscrape} is written in Rust and bundles its templates internally, it can be distributed as a single native executable per target platform. Operators can copy the executable to a restricted node and parse CLI outputs without installing a Python runtime or template package set.

\subsection{Embedded Template Library \& Hot-Reloading}
To eliminate the friction of distributing raw template files, \texttt{cliscrape} ships with a small embedded template library. Using an XDG-compliant resolver, users can parse outputs by declaring a template identifier, e.g., \texttt{-t cisco\_ios\_show\_version}, or an explicit template file. User-defined overrides placed in \texttt{\textasciitilde{}/.local/share/cliscrape/templates/} take precedence over embedded defaults.

Furthermore, the embedded library utilizes \texttt{rust-embed} to provide dual-mode execution. In release builds, templates are compiled directly into the binary for self-contained template deployment. In debug builds, the engine dynamically reads from the local filesystem, enabling instant hot-reloading of templates during development without requiring recompilation.

\subsection{Validation and Edge Case Management}
Network CLI data is notoriously inconsistent due to firmware patches, unexpected banners, and hardware variants. The engine implements strict thresholds and boundaries:
\begin{itemize}
    \item \textbf{Coverage Validation:} A configurable warning system triggers if a parsed output matches fewer fields than a minimum acceptable threshold (e.g., 80\%), acting as a canary for firmware formatting changes.
    \item \textbf{Timeout Guard:} Rust's \texttt{regex} crate avoids backtracking-style regex explosions, and parse calls can also set an elapsed-time timeout guard for operational control.
    \item \textbf{Snapshot Testing:} A comprehensive internal suite relies on \texttt{insta} snapshots, guarding parsed \ac{JSON} output against unintended fixture regressions across updates.
    \item \textbf{Execution Modes:} To balance resilience with strictness, the parser supports a Partial-Match mode (gracefully returning successfully parsed data alongside warnings for missing fields) and a \texttt{--strict} Fail-Fast mode that aborts execution upon the first error.
\end{itemize}

\subsection{Structured Observability}
For production deployments, raw text logs are insufficient. The engine utilises Rust's \texttt{tracing} framework and supports JSON log formatting, capturing module-level lifecycle events, resolution paths, and parsing failures for integration into observability pipelines.

%% =====================================================================
\section{Isomorphic Simulation and Synthetic Generation}
\label{sec:simulation}
%% =====================================================================

One architectural direction in \texttt{cliscrape} is round-trip state machine execution. Traditional parsers represent a unidirectional projection from unstructured text to structured data; \texttt{cliscrape} adds best-effort generation from structured records back into CLI-like text.

\subsection{The Inverse Parser (\texorpdfstring{$FSM^{-1}$}{FSM-1})}
Because the modern YAML/TOML templates preserve rule patterns and field metadata, the parsing engine can be run in a reverse direction for some templates. In this mode, \texttt{cliscrape} consumes structured \ac{JSON} records and uses template rules to synthesize CLI-like text. This is useful, but it is not yet a general inverse-regex solver and should not be described as pixel-accurate for every template.

This capability is a foundation for deterministic network-device simulation. By providing a hypothetical device state, operators can generate representative CLI output for downstream automation tests. The richer simulator, session protocol integration, and vendor-authentic formatting remain planned v3.0 work.

\subsection{Mathematical Verification}
Round-trip generation also provides a mechanism for template verification. A template is round-trip stable for a fixture if raw CLI text parses into \ac{JSON}, generates synthetic text, and re-parses into the same \ac{JSON}. Current verification reports mismatches but does not fail the whole suite for every non-injective template, so stricter oracle semantics remain future hardening work.

%% =====================================================================
\section{Interactive TUI Debugger}
\label{sec:tui}
%% =====================================================================

The built-in \ac{TUI} provides an integrated "Live Lab", offering visual, interactive observability into the parser execution, eliminating trial-and-error template development.

Beyond simply debugging a single file, the TUI includes a \textbf{Template Browser}. Users can press \texttt{t} to browse embedded and XDG user templates, select a template (e.g., \texttt{cisco\_ios\_show\_version.yaml}), and load it into the Live Lab alongside their current raw CLI text. The parsing engine re-runs when watched template/input files change, and after saving edits in the template editor.

\begin{figure}[H]
\centering
\adjustbox{max width=\linewidth}{%
\begin{tikzpicture}[skd, node distance=9mm and 12mm]
  \node[skdPhase, minimum width=11.2cm] (hdr) {\textbf{Debugger control loop: step, breakpoints, run-to-record}};

  \node[skdControl, below left=of hdr, minimum width=3.6cm] (mode) {\textbf{Mode}\\\footnotesize step / run};
  \node[skdControl, right=of mode, minimum width=3.6cm] (bp) {\textbf{Breakpoints}\\\footnotesize state\,=\,S\\line\,=\,i};
  \node[skdCompute, right=of bp, minimum width=3.6cm] (tick) {\textbf{Tick}\\\footnotesize execute 1 match\\or 1 line};

  \node[skdCompute, below=of bp, minimum width=11.2cm] (check) {\textbf{Stop conditions}\\\footnotesize hit breakpoint\quad next Record action\quad end of input};
  \node[skdStore, below=of check, minimum width=11.2cm] (view) {\textbf{UI updates}\\\footnotesize trace append, highlight, vars, diff view};

  \draw[skdCtrlArrow] (mode) -- (tick);
  \draw[skdCtrlArrow] (bp) -- (tick);
  \draw[skdDataArrow] (tick) -- (check);
  \draw[skdDataArrow] (check) -- (view);
  \draw[skdCtrlArrow] (view.west) -- ++(-0.6,0) |- node[near start, left, font=\scriptsize] {user input} (mode.south);
\end{tikzpicture}%
}
\caption{Debugger as a control loop over the parsing interpreter.}
\label{fig:cs-debug-loop}
\end{figure}

The interface offers four panes: an \textbf{Input Pane} for raw text with line highlighting, a \textbf{State Pane} detailing the \ac{FSM} transition history, a \textbf{Variables Pane} showing live value states, and a \textbf{Diff View} explicitly highlighting matching capture groups.


%% =====================================================================
\section{Future Trajectories: The Network Compiler}
\label{sec:future}
%% =====================================================================

By combining a TextFSM-compatible runtime, typed modern templates, semantic diffing, policy constraints, and best-effort generation, \texttt{cliscrape} provides a foundation for several advanced applications within the broader automation ecosystem.

\subsection{Semantic Drift Analysis}
Traditional diff utilities are line-oriented and text-blind, heavily penalizing syntactic noise such as timestamp updates or pagination artifacts. Because \texttt{cliscrape} projects raw text into structured \ac{JSON} records, it enables Operational Diffs. By isolating the delta between two parsed states ($S_1 \rightarrow S_2$), the system identifies semantic transitions while ignoring fields explicitly marked as diff noise.

\subsection{Vendor-Neutral State Schemas}
Network data is traditionally trapped within vendor-specific syntax constraints. \texttt{cliscrape}'s modern YAML definitions include a \texttt{common\_schema} field hook that can map vendor-specific names onto a universal semantic schema. The embedded templates do not yet provide full ledger coverage; standard schemas and compliance validation are the Phase 17 Universal Ledger work.

\subsection{Grammar Induction and Structural Inference}
Current template creation still relies on human refinement. The existing \texttt{infer} command can identify static lines shared across samples and emit broad candidate capture groups as a scaffold. Future iterations should deepen that grammar-induction layer over larger corpuses of raw CLI outputs.

\subsection{The FSM-Oracle: Self-Verification}
The generator facilitates a closed-loop FSM-Oracle direction. When a network device receives a firmware update that alters its CLI output formatting, round-trip checks can detect semantic drift between original and generated records. Today this is a warning-oriented verification loop; a strict CI gate over selected templates is still needed before calling it a proof system.

\begin{figure}[H]
\centering
\adjustbox{max width=\linewidth}{%
\begin{tikzpicture}[skd, node distance=12mm and 15mm]
  \node[skdData, minimum width=3.5cm] (raw) {\textbf{Raw CLI Output}\\\footnotesize (e.g., Firmware v15.4)};
  \node[skdCompute, right=of raw, minimum width=3cm] (parse) {\textbf{Parser}\\\footnotesize $FSM$};
  \node[skdStore, right=of parse, minimum width=3cm] (json) {\textbf{Structured Records}\\\footnotesize JSON};
  \node[skdCompute, below=of parse, minimum width=3cm] (gen) {\textbf{Generator}\\\footnotesize best effort};
  \node[skdData, left=of gen, minimum width=3.5cm] (synth) {\textbf{Synthetic Output}\\\footnotesize generated text};

  \draw[skdDataArrow] (raw) -- (parse);
  \draw[skdDataArrow] (parse) -- (json);
  \draw[skdDataArrow] (json) |- (gen);
  \draw[skdDataArrow] (gen) -- (synth);
  \draw[skdCtrlArrow, <->] (raw.south) -- node[right, font=\scriptsize] {Semantic Round-Trip Check} (synth.north);
\end{tikzpicture}%
}
\caption{The FSM-Oracle direction. A template is stronger when generated synthetic output re-parses to the same semantic records as the original raw CLI output.}
\label{fig:cs-fsm-oracle}
\end{figure}


%% =====================================================================
\section{Performance Evaluation}
\label{sec:performance}
%% =====================================================================

The transition from interpreted Python execution to Rust should yield significant performance gains, and the project includes Criterion benchmarks for throughput tracking. Earlier benchmark runs reported approximately \textbf{4.1 million lines per second}; current release notes should regenerate this number from the checked-in benchmark suite before making public performance claims.

The compiled Rust architecture and DFA-based regex matching avoid Python GIL bottlenecks and backtracking-style regex explosions. A dedicated performance-regression CI gate and refreshed benchmark baselines remain open follow-up work.

%% =====================================================================
\section{Design Summary}
\label{sec:summary}
%% =====================================================================

\begin{table}[H]
\centering
\caption{Key design decisions.}
\label{tab:decisions}
\begin{tabularx}{\linewidth}{lX}
  \toprule
  \textbf{Decision} & \textbf{Rationale} \\
  \midrule
  Rust regex crate & DFA-based $O(n)$ matching; no backtracking vulnerabilities \\
  TextFSM compatibility & Migration path for legacy TextFSM templates \\
  Stateful modifiers & Filldown/List support native relational data extraction \\
  Dual template formats & Incremental adoption; YAML for new templates \\
  Native type coercion & YAML template hints yield directly queryable JSON types \\
  Best-effort generation & Structured records can generate representative synthetic CLI text \\
  Embedded library & XDG-compliant, zero-configuration parsing out-of-the-box \\
  TUI debugger & Interactive FSM visualisation reduces trial-and-error template development \\
  \bottomrule
\end{tabularx}
\end{table}

%% =====================================================================
\printbibliography
%% =====================================================================

\clearpage

% BEGIN AUTOUPDATE
% END AUTOUPDATE

\printindex

\printrevisionhistory
\printtodolist


\end{document}
